\section{On-Device Cell Detection}
\label{sec:edgeai}

\subsection{Inference target}
On-device inference keeps images on the instrument and enables near-real-time
analysis inside the incubator (NFR3). The target accelerator is the Hailo-8 NPU
(\SI{26}{TOPS}) on the AI~HAT+; models must be compiled to the Hailo Executable
Format (HEF) with the Hailo Dataflow Compiler and quantised to integer
precision to run on the device.

\subsection{Model choice}
The biological readout---confluence, count, morphology or migration---fixes both
the model family and the CVAT label schema, and is an open decision (milestone
M1, Section~\ref{sec:requirements}). Bounding-box detection (YOLOv8 and
variants) suffices for counting and confluence, whereas instance segmentation
(Cellpose~\cite{Stringer2021Cellpose}, StarDist~\cite{Schmidt2018StarDist})
is required for per-cell morphology. Only a foreign YOLOv8s-Pose benchmark
currently exists in the project; the task-specific model is future work.
\todo{fix the readout and label schema (M1), then train and compile the
task-specific detector to HEF; report accuracy and on-device latency.}

\subsection{Annotation pipeline}
Training data are annotated in a CVAT server deployed on a separate VPS
(Section~\ref{sec:software}). Acquired chip images are uploaded, labelled
according to the chosen schema, exported and used to train the detector, which
is then quantised and compiled for the Hailo-8. \todo{describe the dataset size,
train/val split and augmentation once acquisition on real chips has started.}

\subsection{Evaluation plan}
The detector will be evaluated both offline (detection metrics on a held-out
set) and on-device (end-to-end latency and throughput per chip on the Pi~5 +
Hailo-8), so that the imaging schedule can be sized against the per-chip
inference time. Results are reported in Section~\ref{sec:evaluation} once the
model exists.
